\documentclass[11pt,a4paper]{article}

% Obsidian to LaTeX Compiler - Preamble
% Adapted from Gilles Castel's lecture notes style

% Core packages
\usepackage[utf8]{inputenc}
\usepackage[T1]{fontenc}
\usepackage{lmodern}

% Math packages
\usepackage{amsmath,amssymb}
\usepackage{mathtools}
\usepackage{bm}

% Layout and formatting
\usepackage{geometry}
\usepackage{parskip}
\usepackage{enumitem}
\usepackage{booktabs}
\usepackage{array}

% Graphics and colors
\usepackage{graphicx}
\usepackage{xcolor}
\usepackage{tikz}
\usetikzlibrary{arrows,positioning,shapes,calc}

% Hyperlinks
\usepackage{hyperref}
\usepackage{url}

% Code highlighting
\usepackage{minted}

% Boxes for callouts
\usepackage{tcolorbox}
\tcbuselibrary{theorems,skins,breakable}

% Quotes
\usepackage{csquotes}

% Text formatting
\usepackage{soul}  % For highlighting (provides \hl)
\usepackage{ulem}  % For strikethrough
\normalem  % Prevent ulem from changing \emph

% Page geometry
\geometry{
    a4paper,
    margin=1in,
    headheight=14pt
}

% Hyperref setup
\hypersetup{
    colorlinks=true,
    linkcolor=blue!70!black,
    filecolor=magenta,
    urlcolor=cyan!70!black,
    citecolor=green!50!black,
    pdftitle={},
    pdfauthor={},
}

% ============================================
% Color definitions for theorem-like boxes
% ============================================

% Primary colors
\definecolor{notecolor}{RGB}{66, 133, 244}       % Blue
\definecolor{examplecolor}{RGB}{52, 168, 83}     % Green
\definecolor{warningcolor}{RGB}{234, 67, 53}     % Red
\definecolor{successcolor}{RGB}{52, 168, 83}     % Green
\definecolor{infocolor}{RGB}{66, 133, 244}       % Blue
\definecolor{abstractcolor}{RGB}{156, 39, 176}   % Purple
\definecolor{tipcolor}{RGB}{0, 150, 136}         % Teal
\definecolor{questioncolor}{RGB}{255, 152, 0}    % Orange
\definecolor{definitioncolor}{RGB}{63, 81, 181}  % Indigo
\definecolor{theoremcolor}{RGB}{103, 58, 183}    % Deep Purple
\definecolor{proofcolor}{RGB}{117, 117, 117}     % Gray
\definecolor{quotecolor}{RGB}{158, 158, 158}     % Light Gray

% ============================================
% Theorem-like environments using tcolorbox
% Note: Using "box" suffix to avoid conflicts with built-in LaTeX environments
% ============================================

% Note box
\newtcolorbox{notebox}[1][Note]{
    enhanced,
    colback=notecolor!5,
    colframe=notecolor!80,
    fonttitle=\bfseries,
    title={#1},
    attach boxed title to top left={yshift=-2mm,xshift=5mm},
    boxed title style={colback=notecolor!80,sharp corners},
    sharp corners,
    breakable
}

% Example box
\newtcolorbox{examplebox}[1][Example]{
    enhanced,
    colback=examplecolor!5,
    colframe=examplecolor!80,
    fonttitle=\bfseries,
    title={#1},
    attach boxed title to top left={yshift=-2mm,xshift=5mm},
    boxed title style={colback=examplecolor!80,sharp corners},
    sharp corners,
    breakable
}

% Warning box
\newtcolorbox{warningbox}[1][Warning]{
    enhanced,
    colback=warningcolor!5,
    colframe=warningcolor!80,
    fonttitle=\bfseries,
    title={#1},
    attach boxed title to top left={yshift=-2mm,xshift=5mm},
    boxed title style={colback=warningcolor!80,sharp corners},
    sharp corners,
    breakable
}

% Success box
\newtcolorbox{successbox}[1][Success]{
    enhanced,
    colback=successcolor!5,
    colframe=successcolor!80,
    fonttitle=\bfseries,
    title={#1},
    attach boxed title to top left={yshift=-2mm,xshift=5mm},
    boxed title style={colback=successcolor!80,sharp corners},
    sharp corners,
    breakable
}

% Info box
\newtcolorbox{infobox}[1][Info]{
    enhanced,
    colback=infocolor!5,
    colframe=infocolor!80,
    fonttitle=\bfseries,
    title={#1},
    attach boxed title to top left={yshift=-2mm,xshift=5mm},
    boxed title style={colback=infocolor!80,sharp corners},
    sharp corners,
    breakable
}

% Abstract box
\newtcolorbox{abstractbox}[1][Abstract]{
    enhanced,
    colback=abstractcolor!5,
    colframe=abstractcolor!80,
    fonttitle=\bfseries,
    title={#1},
    attach boxed title to top left={yshift=-2mm,xshift=5mm},
    boxed title style={colback=abstractcolor!80,sharp corners},
    sharp corners,
    breakable
}

% Tip box
\newtcolorbox{tipbox}[1][Tip]{
    enhanced,
    colback=tipcolor!5,
    colframe=tipcolor!80,
    fonttitle=\bfseries,
    title={#1},
    attach boxed title to top left={yshift=-2mm,xshift=5mm},
    boxed title style={colback=tipcolor!80,sharp corners},
    sharp corners,
    breakable
}

% Question box
\newtcolorbox{questionbox}[1][Question]{
    enhanced,
    colback=questioncolor!5,
    colframe=questioncolor!80,
    fonttitle=\bfseries,
    title={#1},
    attach boxed title to top left={yshift=-2mm,xshift=5mm},
    boxed title style={colback=questioncolor!80,sharp corners},
    sharp corners,
    breakable
}

% Definition box
\newtcolorbox{definitionbox}[1][Definition]{
    enhanced,
    colback=definitioncolor!5,
    colframe=definitioncolor!80,
    fonttitle=\bfseries,
    title={#1},
    attach boxed title to top left={yshift=-2mm,xshift=5mm},
    boxed title style={colback=definitioncolor!80,sharp corners},
    sharp corners,
    breakable
}

% Theorem box
\newtcolorbox{theorembox}[1][Theorem]{
    enhanced,
    colback=theoremcolor!5,
    colframe=theoremcolor!80,
    fonttitle=\bfseries,
    title={#1},
    attach boxed title to top left={yshift=-2mm,xshift=5mm},
    boxed title style={colback=theoremcolor!80,sharp corners},
    sharp corners,
    breakable
}

% Lemma box
\newtcolorbox{lemmabox}[1][Lemma]{
    enhanced,
    colback=theoremcolor!5,
    colframe=theoremcolor!60,
    fonttitle=\bfseries,
    title={#1},
    attach boxed title to top left={yshift=-2mm,xshift=5mm},
    boxed title style={colback=theoremcolor!60,sharp corners},
    sharp corners,
    breakable
}

% Corollary box
\newtcolorbox{corollarybox}[1][Corollary]{
    enhanced,
    colback=theoremcolor!5,
    colframe=theoremcolor!40,
    fonttitle=\bfseries,
    title={#1},
    attach boxed title to top left={yshift=-2mm,xshift=5mm},
    boxed title style={colback=theoremcolor!40,sharp corners},
    sharp corners,
    breakable
}

% Proof box
\newtcolorbox{proofbox}[1][Proof]{
    enhanced,
    colback=proofcolor!5,
    colframe=proofcolor!50,
    fonttitle=\bfseries,
    title={#1},
    attach boxed title to top left={yshift=-2mm,xshift=5mm},
    boxed title style={colback=proofcolor!50,sharp corners},
    sharp corners,
    breakable
}

% Quote box
\newtcolorbox{quotebox}[1][]{
    enhanced,
    colback=quotecolor!5,
    colframe=quotecolor!40,
    left=10pt,
    right=10pt,
    top=5pt,
    bottom=5pt,
    sharp corners,
    breakable
}

% ============================================
% Minted (code highlighting) settings
% ============================================

\setminted{
    frame=lines,
    framesep=2mm,
    linenos=true,
    fontsize=\small,
    breaklines=true,
    tabsize=4,
    bgcolor=gray!5
}

% ============================================
% Custom commands
% ============================================

% Better looking vectors
\renewcommand{\vec}[1]{\mathbf{#1}}

% Common math operators
\DeclareMathOperator{\tr}{tr}
\DeclareMathOperator{\diag}{diag}
\DeclareMathOperator{\rank}{rank}
\DeclareMathOperator{\Span}{span}




\title{Fixed Point Classification}
\author{[]}
\date{2026-01-15T10:25:05-05:00}

\begin{document}

\maketitle

Resuming work on Linear Systems\footnote{\url{https://ewan.my/notes/chaos-theory/linear-systems}} regarding the Logistic Map\footnote{\url{https://ewan.my/notes/selfstudy-mathematics/logistic-map}} using a Cobweb Plot\footnote{\url{https://ewan.my/notes/selfstudy-mathematics/cobweb-plot}}:

A \textbf{fixed point} is a point $p$ such that $f(p) = p$. Usually a sink/source, but due to periodic points, can be a rotation set of points.

\begin{displayquote}
The \textbf{epsilon-neighborhood} $N_{\epsilon}(p)$ of a point $p$ is the internal $\{ x \in \mathbb{R}: |x-p| < \epsilon \}$.
\end{displayquote}

\begin{displayquote}
If $\exists \epsilon > 0 : \forall x  \in N_{\epsilon}(p)$, $\lim_{ k \to \infty }f^{k}(x) = p$ then we call $p$ a \textbf{sink} or \textbf{Attracting fixed point\footnote{\url{https://ewan.my/notes/chaos-theory/attracting-fixed-point}}}
\end{displayquote}

In plain English, if there is some neighborhood around point $p$ such that all points in that neighborhood converge to $p$ under iteration of $f$, then $p$ is a sink.

\begin{displayquote}
If $\exists \epsilon > 0 :$ all $x \in N_{\epsilon}(p)$ except $p$ eventually map outside of $N_{\epsilon}(p)$, then we call $p$ a \textbf{source} or a \textbf{Repelling fixed point\footnote{\url{https://ewan.my/notes/chaos-theory/repelling-fixed-point}}}.
\end{displayquote}

In plain English, if there is some neighborhood around point $p$ such that all points in that neighborhood (except for $p$ itself) eventually leave that neighborhood under iteration of $f$, then $p$ is a source.

_For cases in this class, $\epsilon$ represents a small positive real number parameter._

\hrulefill

\subsection{Cubic Example}

Looking at the function:

\[
f(x) = \frac{3x-x^3}{2}
\]

\begin{infobox}[Desmos]

<iframe src="https://www.desmos.com/calculator/mqr1jsj7l9" width=600 height="400" ></iframe>
\end{infobox}

The \textbf{fixed point sink\footnote{\url{https://ewan.my/notes/chaos-theory/attracting-fixed-point}}s are}: $x = \pm 1$

The \textbf{fixed point Source\footnote{\url{https://ewan.my/notes/chaos-theory/repelling-fixed-point}}s are}: $x=0$

\textbf{Key observations:}

\begin{itemize}
  \item Sinks at $x = \pm 1$ are \textbf{super-attracting}: $f'(\pm 1) = 0$ (critical points = fixed points)
  \item Source at $x = 0$ has $f'(0) = \frac{3}{2} > 1$ (linearly unstable)
  \item \textbf{Odd symmetry}: $f(-x) = -f(x)$ creates symmetric basins of attraction
  \item Topologically similar to Logistic Map\footnote{\url{https://ewan.my/notes/selfstudy-mathematics/logistic-map}} at $r=3$ (pre-chaotic regime)
\end{itemize}

\begin{abstractbox}[Theorem]

Let $f$ be a smooth map (derivatives continousous), and assume $p$ is fixed point of $f$.

If $|f'(p)| <1 \Rightarrow p$ is a sink\footnote{\url{https://ewan.my/notes/chaos-theory/attracting-fixed-point}}.

If $|f'(p)| > 1 \Rightarrow p$ is a Source\footnote{\url{https://ewan.my/notes/chaos-theory/repelling-fixed-point}}.

If $|f'(p)| =1$, we need more information.
\end{abstractbox}

Origin is a "neutral" fixed point, all other $x \in \mathbb{R}$ are period 2. In essence, some starting position will take a minute to read the period two orbit, whereas others like $p_{1}$ will already be on the period two orbit.

\hrulefill

\subsection{Periodic Sinks and Source}

\begin{displayquote}
Let $f$ be a map and assume $p$ is a period $k$ point. The period k orbit of $p$ is a periodic sink (source) if $p$ is a sink( source ) for the map $f^{k}$
\end{displayquote}

In plain English, if $p$ is a period $k$ point, then we can look at the $k$th iterate of $f$ and classify $p$ as a sink or source based on that.




\end{document}
